% \begin{document}
\chapter{Logical Relations}
 
\section{Introduction}
 
The simply-typed $\lambda$-calculus rules out the nonterminating
$\Omega = (\lambda x.\, x\,x)\,(\lambda x.\, x\,x)$: typing $x\,x$ would need a type
$\tau$ with $\tau = \tau \to \tau'$, and none exists. In fact every well-typed program
terminates.
 
\thm{Termination}{
If $\vdash e : \tau$, then there is a value $v$ such that $e \to^* v$.
}
 
We prove this using \emph{logical relations}, a technique invented by Tait in 1967.
Termination is a strong guarantee, on which total languages such as Rocq (Coq), Agda,
and Lean rely for logical consistency. Unlike type soundness it is not a syntactic
property: progress and preservation say a well-typed program never gets stuck, but
nothing about how long it runs. The technique matters even more than the theorem;
logical relations are used to prove
\begin{itemize}
  \item type soundness (semantically) and parametricity and free theorems (next lecture),
  \item program equivalence and compiler correctness,
  \item security properties such as noninterference, and
  \item safety of libraries that use \texttt{unsafe} code in Rust.
\end{itemize}
 
\section{The Language}
 
We work with the simply-typed $\lambda$-calculus over $\kw{unit}$:
\begin{align*}
e &::= x \mid \lambda x : \tau.\, e \mid e_1\, e_2 \mid () \\
v &::= \lambda x : \tau.\, e \mid () \\
\tau &::= \kw{unit} \mid \tau_1 \to \tau_2
\end{align*}
with standard call-by-value semantics
\[
E ::= [\cdot] \mid E\, e \mid v\, E
\qquad
\inferrule*[Right=\textsc{Context}]
  {e \to e'}
  {E[e] \to E[e']}
\qquad \qquad
\inferrule*[Right=\textsc{$\beta$-Reduction}]
  { }
  {(\lambda x : \tau.\, e)\, v \to \tsub{e}{v}{x}}
\]
and the typing rules \textsc{T-Var}, \textsc{T-Unit}, \textsc{T-Abs}, \textsc{T-App}
from Lecture 12.
 
\dfn{Halting}{ We write $e \halts$, pronounced ``$e$
halts'', if $e \to^* v$ for some value $v$.
} 
\noindent We use these standard facts from Lecture 12:
\begin{itemize}
  \item \emph{Determinism}: if $e \to e_1$ and $e \to e_2$ then $e_1 = e_2$. So if
        $e \to e'$, then $e \halts$ if and only if $e' \halts$.
  \item \emph{Preservation}: if $\vdash e : \tau$ and $e \to e'$ then $\vdash e' : \tau$.
  \item \emph{Canonical forms}: if $\vdash v : \kw{unit}$ then $v = ()$; if
        $\vdash v : \tau_1 \to \tau_2$ then $v = \lambda x : \tau_1.\, e$ for some $x$
        and $e$.
  \item \emph{Substitution}: if $\Gamma, x : \tau' \vdash e : \tau$ and
        $\vdash v : \tau'$ then $\Gamma \vdash \tsub{e}{v}{x} : \tau$.
\end{itemize}
 
\section{A First Attempt}
 
We try to prove that $\vdash e : \tau$ implies $e \halts$ by induction on the typing
derivation.
\begin{itemize}
  \item \textsc{T-Var}: impossible, the context is empty.
  \item \textsc{T-Unit}, \textsc{T-Abs}: immediate, since $()$ and $\lambda x : \tau.\, e'$
        are values.
  \item \textsc{T-App}: $e = e_1\, e_2$ with $\vdash e_1 : \tau' \to \tau$ and
        $\vdash e_2 : \tau'$. By the induction hypotheses $e_1 \to^* v_1$ and
        $e_2 \to^* v_2$; by preservation (actually type soundness because of multistep) and canonical forms
        $v_1 = \lambda x : \tau'.\, e'$. So
        \[
        e_1\, e_2 \to^* (\lambda x : \tau'.\, e')\, e_2 \to^* (\lambda x : \tau'.\, e')\, v_2 \to \tsub{e'}{v_2}{x}.
        \]
        We would like $\tsub{e'}{v_2}{x} \halts$, but no induction hypothesis applies:
        it is not a subexpression of $e$, its typing derivation is not part of $e$'s,
        and it may be larger than $e$, so induction on size does not help either.
\end{itemize}
The induction hypothesis for $e_1$ says only that it halts, not what happens when it
is applied to an argument. We need a stronger hypothesis.
 
\section{Logical Relations}
 
For function types we need to know that the function halts \emph{and} that applying it
to a halting argument gives a halting result. For higher-order functions arguments and
results are themselves functions, so the same property is needed for them, and so on.
We therefore define the property by induction on the type; a family of predicates
indexed by types is a \emph{logical relation}, here capturing \emph{hereditary
termination}.
 
\dfn{Logical Relation}{$R_\tau(e)$ is defined by induction
on $\tau$:
\begin{align*}
R_{\kw{unit}}(e) &\iff\ \vdash e : \kw{unit} \ \text{ and } \ e \halts \\
R_{\tau_1 \to \tau_2}(e) &\iff\ \vdash e : \tau_1 \to \tau_2 \ \text{ and } \ e \halts \ \text{ and }\
  \text{for every } e',\ R_{\tau_1}(e') \implies R_{\tau_2}(e\, e')
\end{align*}
}
This is well-founded, since the clause for $\tau_1 \to \tau_2$ uses only $R_{\tau_1}$
and $R_{\tau_2}$ at smaller types. The function clause resembles a logical implication
(related arguments to related results), hence ``logical''.
 
Termination follows from two lemmas. The first says related expressions halt.
 
\mlenma{Adequacy}{
If $R_\tau(e)$, then $e \halts$.
}
 
\pf{Proof}{
By cases on $\tau$. If $\tau = \kw{unit}$, then $R_{\kw{unit}}(e)$ means
$\vdash e : \kw{unit}$ and $e \halts$; in particular $e \halts$. If
$\tau = \tau_1 \to \tau_2$, then $R_{\tau_1 \to \tau_2}(e)$ means
$\vdash e : \tau_1 \to \tau_2$, $e \halts$, and $e$ maps related arguments to related
results; in particular $e \halts$. (Halting is built into every clause, so this is by
cases rather than induction.)
}
 
The second says well-typed expressions are related.
 
\mlenma{Fundamental lemma, first version}{
If $\vdash e : \tau$, then $R_\tau(e)$.
}
 
First we need a lemma about evaluation.
 
\mlenma{Closure under evaluation}{
If $\vdash e : \tau$ and $e \to e'$, then $R_\tau(e)$ if and only if $R_\tau(e')$.
}
 
\pf{Proof}{
By induction on $\tau$.
 
\emph{Case $\tau = \kw{unit}$.} $(\Rightarrow)$ Suppose $R_{\kw{unit}}(e)$. Then
$\vdash e' : \kw{unit}$ by preservation, and since $e \halts$ and $e \to e'$,
determinism gives $e' \halts$. So $R_{\kw{unit}}(e')$. $(\Leftarrow)$ Suppose
$R_{\kw{unit}}(e')$. Then $\vdash e : \kw{unit}$ by assumption, and since $e \to e'$
and $e' \halts$, we have $e \halts$. So $R_{\kw{unit}}(e)$.
 
\emph{Case $\tau = \tau_1 \to \tau_2$.} Let $e''$ be any expression with
$R_{\tau_1}(e'')$. Then $\vdash e'' : \tau_1$, so
\begin{equation}
\vdash e\, e'' : \tau_2 \quad\text{and}\quad e\, e'' \to e'\, e'', \tag{$*$}
\end{equation}
where the step uses the evaluation context $E\, e''$. So the induction hypothesis at
$\tau_2$ applies to $e\, e''$ and $e'\, e''$.
 
$(\Rightarrow)$ Suppose $R_{\tau_1 \to \tau_2}(e)$. We show the parts of
$R_{\tau_1 \to \tau_2}(e')$: $\vdash e' : \tau_1 \to \tau_2$ by preservation;
$e' \halts$ since $e \halts$ and $e \to e'$ (determinism); and for $R_{\tau_1}(e'')$,
since $R_{\tau_1 \to \tau_2}(e)$ we have $R_{\tau_2}(e\, e'')$, so by $(*)$ and the
induction hypothesis at $\tau_2$ (left to right), $R_{\tau_2}(e'\, e'')$.
 
$(\Leftarrow)$ Suppose $R_{\tau_1 \to \tau_2}(e')$. We show the parts of
$R_{\tau_1 \to \tau_2}(e)$: $\vdash e : \tau_1 \to \tau_2$ by assumption; $e \halts$
since $e \to e'$ and $e' \halts$; and for $R_{\tau_1}(e'')$, since
$R_{\tau_1 \to \tau_2}(e')$ we have $R_{\tau_2}(e'\, e'')$, so by $(*)$ and the
induction hypothesis at $\tau_2$ (right to left), $R_{\tau_2}(e\, e'')$.
}
 
By induction on the derivation of $e \to^* e'$ the same holds for the multi-step
relation, used in both directions:
\begin{itemize}
  \item \emph{forward}: if $R_\tau(e)$ and $e \to^* v$, then $R_\tau(v)$;
  \item \emph{backward}: if $\vdash e : \tau$, $e \to^* e'$, and $R_\tau(e')$, then
        $R_\tau(e)$.
\end{itemize}
 
\subsection{The Fundamental Lemma}
 
\pf{Proof attempt of the first version}{
By induction on the derivation of $\vdash e : \tau$. In every case $\vdash e : \tau$
holds by assumption, so we only show the rest of $R_\tau(e)$.
 
\emph{\textsc{T-Var}}: impossible, the context is empty.
 
\emph{\textsc{T-Unit}}: $e = ()$ and $\tau = \kw{unit}$; $()$ is a value, so it halts
and $R_{\kw{unit}}(())$.
 
\emph{\textsc{T-App}}: $e = e_1\, e_2$ with $\vdash e_1 : \tau' \to \tau$ and
$\vdash e_2 : \tau'$. By the induction hypotheses $R_{\tau' \to \tau}(e_1)$ and
$R_{\tau'}(e_2)$. By definition of $R_{\tau' \to \tau}(e_1)$, for every $e'$ with
$R_{\tau'}(e')$ we have $R_\tau(e_1\, e')$; taking $e' = e_2$ gives
$R_\tau(e_1\, e_2)$. Unlike the first attempt, we never look at how $e_1\, e_2$
evaluates: the stronger hypothesis does all the work.
 
\emph{\textsc{T-Abs}}: $e = \lambda x : \tau_1.\, e_0$ and $\tau = \tau_1 \to \tau_2$
with $x : \tau_1 \vdash e_0 : \tau_2$. We must show (1)
$\vdash e : \tau_1 \to \tau_2$, which holds by assumption; (2) $e \halts$, since $e$
is a value; (3) for every $e'$, $R_{\tau_1}(e') \implies R_{\tau_2}(e\, e')$. For (3),
let $R_{\tau_1}(e')$. By adequacy $e' \to^* v'$ for some $v'$, and by closure under
(multi)evaluation (forward) $R_{\tau_1}(v')$. Then
\[
(\lambda x : \tau_1.\, e_0)\, e' \to^* (\lambda x : \tau_1.\, e_0)\, v' \to \tsub{e_0}{v'}{x}.
\]
Since $\vdash e\, e' : \tau_2$, closure under evaluation (backward) reduces the goal
to $R_{\tau_2}(\tsub{e_0}{v'}{x})$.
 
Here we are stuck. The premise of \textsc{T-Abs} is $x : \tau_1 \vdash e_0 : \tau_2$,
so $e_0$ is open, while the induction hypothesis concerns only closed expressions and
says nothing about $e_0$, let alone $\tsub{e_0}{v'}{x}$.
}
 
The solution is to strengthen the lemma to open expressions, closing them with a
substitution of related values for free variables. A substitution
$\gamma = \{v_1/x_1, \ldots, v_k/x_k\}$ maps variables to closed values. We write
$\gamma(e)$ for the result of substituting each $v_i$ for $x_i$ in $e$ (the order is
irrelevant, since the values are closed), and $\gamma[x \mapsto v]$ for the
substitution mapping $x$ to $v$ and otherwise like $\gamma$.
 
\dfn{Related substitutions}{
\[
\gamma \models \Gamma \iff \dom(\gamma) = \dom(\Gamma) \text{ and } R_{\Gamma(x)}(\gamma(x)) \text{ for each } x \in \dom(\Gamma).
\]
} 
\mlenma{Fundamental lemma}{
If $\Gamma \vdash e : \tau$ and $\gamma \models \Gamma$, then $R_\tau(\gamma(e))$.
}
 
With $\Gamma$ and $\gamma$ empty this is the first version.
 
\pf{Proof}{
By induction on the derivation of $\Gamma \vdash e : \tau$. Each $\gamma(x)$ is a
closed value with $\vdash \gamma(x) : \Gamma(x)$, so the substitution lemma gives
$\vdash \gamma(e) : \tau$ in every case; we only show the rest of $R_\tau(\gamma(e))$.
 
\emph{\textsc{T-Var}}: $e = x$ and $\tau = \Gamma(x)$. Since $\gamma \models \Gamma$,
$R_{\Gamma(x)}(\gamma(x))$.
 
\emph{\textsc{T-Unit}}: $e = ()$, $\tau = \kw{unit}$, and $\gamma(e) = ()$ halts.
 
\emph{\textsc{T-App}}: $e = e_1\, e_2$ with $\Gamma \vdash e_1 : \tau' \to \tau$ and
$\Gamma \vdash e_2 : \tau'$. By the induction hypotheses
$R_{\tau' \to \tau}(\gamma(e_1))$ and $R_{\tau'}(\gamma(e_2))$. By definition of
$R_{\tau' \to \tau}$ we get $R_\tau(\gamma(e_1)\, \gamma(e_2))$, and
$\gamma(e_1)\, \gamma(e_2) = \gamma(e_1\, e_2)$.
 
\emph{\textsc{T-Abs}}: $e = \lambda x : \tau_1.\, e_0$ and $\tau = \tau_1 \to \tau_2$
with $\Gamma, x : \tau_1 \vdash e_0 : \tau_2$. By renaming assume $x \notin \dom(\Gamma)$,
so $\gamma(e) = \lambda x : \tau_1.\, \gamma(e_0)$, a value, hence halting. Let
$R_{\tau_1}(e')$.
\begin{itemize}
  \item By adequacy $e' \to^* v'$ for some $v'$, and by closure under evaluation
        $R_{\tau_1}(v')$.
  \item So $\gamma[x \mapsto v'] \models \Gamma, x : \tau_1$, and by the induction
        hypothesis $R_{\tau_2}(\gamma[x \mapsto v'](e_0))$, i.e.\
        $R_{\tau_2}(\tsub{\gamma(e_0)}{v'}{x})$.
  \item But $\gamma(e)\, e' \to^* \gamma(e)\, v' \to \tsub{\gamma(e_0)}{v'}{x}$ and
        $\vdash \gamma(e)\, e' : \tau_2$, so by closure under evaluation
        $R_{\tau_2}(\gamma(e)\, e')$.
\end{itemize}
Therefore $R_{\tau_1 \to \tau_2}(\gamma(e))$.
}
 
Note that in the abstraction case the induction hypothesis is used for \emph{every}
related value $v'$; this is why the lemma must quantify over all related substitutions
$\gamma$.
 
\thm{Termination}{
If $\vdash e : \tau$, then $e \halts$.
}
 
\pf{Proof}{
By the fundamental lemma with $\Gamma$ and $\gamma$ empty, $R_\tau(e)$. By adequacy,
$e \halts$.
}
 
A proof by logical relations follows a recipe:
\begin{enumerate}
  \item Define a relation by induction on types, with the property of interest at base
        types and ``related arguments to related results'' at function types.
  \item Prove \emph{adequacy}: related terms have the property.
  \item Prove the \emph{fundamental lemma}: well-typed terms are related, generalized to
        open terms closed by related substitutions.
\end{enumerate}
 
\section{Extensions}
 
\subsection{Booleans}
 
Add $\ttkw{true}$, $\ttkw{false}$, and $\ttkw{if}\ e_1\ \ttkw{then}\ e_2\ \ttkw{else}\
e_3$ with the semantics and typing rules from Lecture 12. The relation gains one
clause:
\[
R_{\kw{bool}}(e) \iff\ \vdash e : \kw{bool} \text{ and } e \halts.
\]
Closure under evaluation at $\kw{bool}$ is as for $\kw{unit}$. The fundamental lemma
gets three new cases; \textsc{T-True} and \textsc{T-False} are immediate since
constants are values.
 
\emph{\textsc{T-If}}: $e = \ttkw{if}\ e_1\ \ttkw{then}\ e_2\ \ttkw{else}\ e_3$ with
$\Gamma \vdash e_1 : \kw{bool}$, $\Gamma \vdash e_2 : \tau$, $\Gamma \vdash e_3 : \tau$.
By the induction hypothesis $R_{\kw{bool}}(\gamma(e_1))$, so $\gamma(e_1) \to^* v$ for
some $v$, which by preservation and canonical forms is $\ttkw{true}$ or
$\ttkw{false}$. If $v = \ttkw{true}$,
\[
\gamma(e) \to^* \ttkw{if}\ \ttkw{true}\ \ttkw{then}\ \gamma(e_2)\ \ttkw{else}\ \gamma(e_3) \to \gamma(e_2).
\]
By the induction hypothesis $R_\tau(\gamma(e_2))$, so by closure under evaluation
$R_\tau(\gamma(e))$. The case $v = \ttkw{false}$ is symmetric. We never needed to know
anything about the branch not taken.
 
\subsection{Products}
 
An expression of product type is related if it halts and both projections are related:
\[
R_{\tau_1 \times \tau_2}(e) \iff\ \vdash e : \tau_1 \times \tau_2 \text{ and } e \halts \text{ and } R_{\tau_1}(\#1\ e) \text{ and } R_{\tau_2}(\#2\ e).
\]
In general each type constructor gets a clause saying how its values are used:
functions are applied, pairs are projected, Booleans are tested. Sums are similar and
are a homework exercise.
 
\subsection{What breaks}
 
\emph{Fixed points.} Adding $\ttkw{fix}$ (Lecture 12) makes the termination theorem
false, so the proof must fail. It fails in the case for
$\ttkw{fix}\ (\lambda f : \tau.\, e)$, which steps to
$\tsub{e}{\ttkw{fix}\ (\lambda f : \tau.\, e)}{f}$: to use the induction hypothesis for
$e$ we would need $R_\tau$ of the $\ttkw{fix}$ expression itself, which is exactly what
we are trying to prove.
 
\emph{Recursive types.} With the recursive types of Lecture 13 we cannot even define
the relation: the clause for $\mu\alpha.\, \tau$ would refer to
$R_{\tsub{\tau}{\mu\alpha.\,\tau}{\alpha}}$, at a type that is not smaller.
 
\subsection{Strong normalization}
 
We proved that CBV evaluation of a well-typed program terminates. The strong
normalization theorem from Lecture 12 says every sequence of full $\beta$-reductions is
finite. The same proof structure works, with ``$e$ is strongly normalizing'' in place of
``$e \halts$''; closure under reduction is more delicate, since reduction is no longer
deterministic.
 
For System F, defining the relation is harder: the relation for $\forall X.\, \tau$ must
say what happens for every instantiation of $X$, including $\forall X.\, \tau$ itself,
so it cannot be defined by induction on types. Girard (1972) solved this by quantifying
over all sets of terms satisfying certain closure properties, called \emph{reducibility
candidates}. The next lecture gives a similar idea in a binary logical relation for
System F.
% \end{document}
